% ============================================================
%  Section 2 — Data
% ============================================================

\section{Data}
\label{sec:data}

\subsection{Auction transactions}

Our primary dataset consists of lot-level hammer prices from French fine-wine auctions covering
the period 1996--2015. Each observation records the château or domaine, vintage year, bottle
format, auction house, sale date, and realised price per bottle in USD. After applying the
sample selection criteria described below, the working sample comprises 1,068,703 lots.
We restrict attention to three regions: Bordeaux (615,078 lots), Burgundy (328,371 lots), and
the Rhône Valley (105,450 lots).

\subsection{Sample selection}
\label{sec:data:selection}

We apply the following sequential exclusions to arrive at the analysis sample:

\begin{enumerate}
  \item \textbf{Non-positive prices.} We drop lots with a hammer price at or below zero
        (data-entry errors and passed lots recorded as zero). \textit{[Report count dropped.]}
  \item \textbf{Non-standard bottle formats.} We retain only standard (750\,ml) and
        magnum (1.5\,L) formats; jeroboams and larger are excluded because their price-per-bottle
        series is thin and driven by different demand. \textit{[Report count dropped.]}
  \item \textbf{Mixed lots.} Lots containing more than one châteu or vintage within the same
        hammer price are flagged and excluded, as the per-bottle price cannot be attributed to a
        single wine. \textit{[Report count dropped.]}
  \item \textbf{Extreme price outliers.} \textit{[Describe outlier trimming rule, if any —
        e.g., top/bottom 0.1\% of log prices within region-vintage cells.]}
  \item \textbf{Vintage range.} We restrict to vintages 1950--2014 to ensure at least one
        year of auction history and exclude pre-Phylloxera vintages for which maturity data
        are unavailable. \textit{[Confirm exact vintage cutoffs.]}
\end{enumerate}

\noindent Table~\ref{tab:full_sample} in Appendix~\ref{app:full_sample} reports lot counts
and median prices for the full pre-exclusion sample broken down by region and wine type.
\textit{[Add an attrition table showing the count at each exclusion step.]}

\subsection{Maturity windows from Tastingbook}
\label{sec:data:tastingbook}

A central feature of our approach is that maturity dates vary across wines: a top Bordeaux
Cru Classé may peak at 25--30 years after vintage, while a lighter Burgundy village wine may
peak at 8--12 years. Treating all wines as maturing on the same calendar would conflate wines
at very different points in their life-cycles. We therefore match each auction lot to a
wine-specific maturity window.

We match auction lots to drink windows from Tastingbook, a comprehensive online
cellar-management database maintained by professional sommeliers and oenologists.
For each château--vintage pair, Tastingbook records a ``drink-from'' year and a
``drink-until'' year; we define the maturity peak $\tau_i$ as the midpoint of this window,
measured in years since the vintage.
\textit{[Add a sentence or two about the scraping procedure and coverage — how many
châteaux-vintages are in Tastingbook, how we matched to auction records (on château name
and vintage year), and any name-standardisation steps.]}

We match 75.9\,\% of auction lots to a Tastingbook record. For the remaining 24.1\,\%, we
impute $\tau_i$ using a gradient-boosting model trained on the matched observations
(cross-validated $R^2 = 0.991$; see Appendix~\ref{app:imputation}). Features for
imputation include the producer, appellation, region, grape variety, and vintage year.

All maturity values are capped at 20 years. Tastingbook recommendations for old vintages
sometimes anchor to the current date (``drink now''), inflating years-to-maturity to
implausible values; 20 years is a conservative and defensible upper bound given the
sample period and the distribution of grand cru maturities.
\textit{[Consider adding a figure showing the distribution of $\tau_i$ across wine types,
illustrating the cross-sectional variation that powers our identification.]}

\subsection{Normalised age}

We normalise wine age by the wine-specific maturity calendar. Define
\begin{equation}
  a_{it}^* = \frac{t - v_i}{\tau_i},
  \label{eq:age_norm}
\end{equation}
where $t$ is the sale year, $v_i$ is the vintage year, and $\tau_i$ is the maturity-peak
year. By construction $a_{it}^* = 1$ when the wine is sold exactly at its maturity peak;
$a_{it}^* < 1$ indicates pre-maturity and $a_{it}^* > 1$ indicates post-maturity. Across
all lots, 72.3\,\% are sold pre-maturity, consistent with speculative and investment demand
dominating auction supply.

\subsection{Grand cru subsamples}

Our main analysis focuses on the highest-quality wines, where the consumer/collector distinction
is sharpest. We define three groups:

\begin{enumerate}
  \item \textbf{Bordeaux Crus Classés (red):} all classified growths from the Médoc, Graves,
        and Saint-Émilion (436,413 lots).
  \item \textbf{Burgundy Grand Cru:} all Grand Cru appellations from the Côte d'Or (221,512 lots).
  \item \textbf{Burgundy Premier Cru:} all Premier Cru appellations from the Côte d'Or
        (87,136 lots).
\end{enumerate}

Summary statistics for these groups are reported in Table~\ref{tab:summary_stats}.
